\documentclass{beamer}\usepackage[]{graphicx}\usepackage[]{color}
%% maxwidth is the original width if it is less than linewidth
%% otherwise use linewidth (to make sure the graphics do not exceed the margin)
\makeatletter
\def\maxwidth{ %
  \ifdim\Gin@nat@width>\linewidth
    \linewidth
  \else
    \Gin@nat@width
  \fi
}
\makeatother

\definecolor{fgcolor}{rgb}{0.345, 0.345, 0.345}
\newcommand{\hlnum}[1]{\textcolor[rgb]{0.686,0.059,0.569}{#1}}%
\newcommand{\hlstr}[1]{\textcolor[rgb]{0.192,0.494,0.8}{#1}}%
\newcommand{\hlcom}[1]{\textcolor[rgb]{0.678,0.584,0.686}{\textit{#1}}}%
\newcommand{\hlopt}[1]{\textcolor[rgb]{0,0,0}{#1}}%
\newcommand{\hlstd}[1]{\textcolor[rgb]{0.345,0.345,0.345}{#1}}%
\newcommand{\hlkwa}[1]{\textcolor[rgb]{0.161,0.373,0.58}{\textbf{#1}}}%
\newcommand{\hlkwb}[1]{\textcolor[rgb]{0.69,0.353,0.396}{#1}}%
\newcommand{\hlkwc}[1]{\textcolor[rgb]{0.333,0.667,0.333}{#1}}%
\newcommand{\hlkwd}[1]{\textcolor[rgb]{0.737,0.353,0.396}{\textbf{#1}}}%

\usepackage{framed}
\makeatletter
\newenvironment{kframe}{%
 \def\at@end@of@kframe{}%
 \ifinner\ifhmode%
  \def\at@end@of@kframe{\end{minipage}}%
  \begin{minipage}{\columnwidth}%
 \fi\fi%
 \def\FrameCommand##1{\hskip\@totalleftmargin \hskip-\fboxsep
 \colorbox{shadecolor}{##1}\hskip-\fboxsep
     % There is no \\@totalrightmargin, so:
     \hskip-\linewidth \hskip-\@totalleftmargin \hskip\columnwidth}%
 \MakeFramed {\advance\hsize-\width
   \@totalleftmargin\z@ \linewidth\hsize
   \@setminipage}}%
 {\par\unskip\endMakeFramed%
 \at@end@of@kframe}
\makeatother

\definecolor{shadecolor}{rgb}{.97, .97, .97}
\definecolor{messagecolor}{rgb}{0, 0, 0}
\definecolor{warningcolor}{rgb}{1, 0, 1}
\definecolor{errorcolor}{rgb}{1, 0, 0}
\newenvironment{knitrout}{}{} % an empty environment to be redefined in TeX

\usepackage{alltt} 
% \usepackage{graphicx}
\usepackage{graphics}
\usepackage[T1]{fontenc}
\usepackage{hyperref}
\usepackage{verbatim}
\setbeamercovered{transparent}
\renewcommand{\ni}{\noindent}
\hypersetup{
  colorlinks   = true, %Colours links instead of ugly boxes
  urlcolor     = blue, %Colour for external hyperlinks
  linkcolor    = blue, %Colour of internal links
  citecolor   = red %Colour of citations
}
%% to include page numbers manually include the next three lines
% \usepackage{fancyhdr,lastpage}
% \pagestyle{fancy}\fancyhf{}\rfoot{\vspace{-0.5cm} Page {\thepage} of \pageref{LastPage}}
% \renewcommand\headrulewidth{0pt} % Removes funny header line
%load packages that will be invisible on slides





\title[Dynamic Documents with knitr]{LaTeX packages for R and Advanced knitr}
\date{April 9, 2014}
\institute[ISU]{Iowa State University}
\IfFileExists{upquote.sty}{\usepackage{upquote}}{}
\begin{document}
%---------------------------------------------------------------------------

\begin{frame}
    \maketitle
\end{frame}

%---------------------------------------------------------------------------

\begin{frame}[fragile]
\frametitle{More ways to combine R and LaTeX}
\begin{itemize}
\item Additional knitr options for formatting R output: \texttt{\textbackslash Sexpr\{\}}, \texttt{results='asis'}
\item \texttt{xtable} - formats R tables and data frames as LaTeX tables\medskip
\item \texttt{stargazer} - displaying R models in LaTeX\medskip
\item \texttt{Hmisc} - display generic R objects in LaTeX\medskip
\end{itemize}
\end{frame}

%---------------------------------------------------------------------------

\begin{frame}[fragile]
\frametitle{Referencing R objects in the LaTeX document}
\begin{knitrout}\footnotesize
\definecolor{shadecolor}{rgb}{1, 1, 1}\color{fgcolor}\begin{kframe}
\begin{alltt}
\hlstd{x} \hlkwb{<-} \hlnum{15}
\hlstd{y} \hlkwb{<-} \hlkwd{rnorm}\hlstd{(}\hlnum{10}\hlstd{,} \hlnum{1}\hlstd{)}
\end{alltt}
\end{kframe}
\end{knitrout}

Often, we want to reference numerical results in the text: reporting summary statistics, number of observations, p-values... \\\hfill we don't want to have to replace those manually.\bigskip\bigskip

\texttt{\textbackslash Sexpr\{R code\}} lets you reference your R data inline. 
\end{frame}

%---------------------------------------------------------------------------

\begin{frame}[fragile]
\frametitle{Referencing R objects in the LaTeX document}
\texttt{\textbackslash Sexpr\{R code\}} lets you reference your R data inline. \bigskip\bigskip

For example,
\begin{knitrout}\footnotesize
\definecolor{shadecolor}{rgb}{1, 1, 1}\color{fgcolor}\begin{kframe}
\begin{verbatim}
$x=\Sexpr{x}$ and $\overline{y}=\Sexpr{mean(y)}$
\end{verbatim}
\end{kframe}
\end{knitrout}

produces the output $x=15$ and $\overline{y}=0.7226$ when compiled in a knitr document. As long as the code chunk you are referencing \textbf{precedes} the \texttt{\textbackslash Sexpr\{\}} command, knitr will be able to fill in the blanks for you!

\end{frame}

%---------------------------------------------------------------------------

\begin{frame}[fragile]
\frametitle{Creating LaTeX code within R}
R has lots of packages to produce LaTeX formatted R objects: 
\begin{itemize}
\item \href{http://cran.r-project.org/web/packages/xtable/index.html}{xtable} - make nice tables
\item \href{http://cran.r-project.org/web/packages/stargazer/index.html}{stargazer} - regression model tables
\item \href{http://cran.r-project.org/web/packages/Hmisc/index.html}{Hmisc} - make generic R objects into LaTeX-formatted objects\\
This one is a bit higher-level
\item \href{http://cran.r-project.org/web/packages/texreg/index.html}{texreg} - model output
\item \href{http://cran.r-project.org/web/packages/miscFuncs/index.html}{miscFuncs} - more latex tables
\item \href{http://cran.r-project.org/web/packages/reporttools/index.html}{reporttools} - descriptive statistics
\end{itemize}

We're only going to talk about the first 3, but there are many other packages out there to do similar things.
\end{frame}

%---------------------------------------------------------------------------

\begin{frame}[fragile]
\frametitle{Creating LaTeX Tables with xtable}
\begin{knitrout}\footnotesize
\definecolor{shadecolor}{rgb}{1, 1, 1}\color{fgcolor}\begin{kframe}
\begin{alltt}
\hlkwd{library}\hlstd{(xtable)}
\hlkwd{data}\hlstd{(iris)}
\hlkwd{xtable}\hlstd{(}\hlkwd{head}\hlstd{(iris))}
\end{alltt}
\end{kframe}
\end{knitrout}

{\footnotesize
% latex table generated in R 3.0.2 by xtable 1.7-1 package
% Wed Mar  5 15:14:09 2014
\begin{table}[ht]
\centering
\begin{tabular}{rrrrrl}
  \hline
 & Sepal.Length & Sepal.Width & Petal.Length & Petal.Width & Species \\ 
  \hline
1 & 5.10 & 3.50 & 1.40 & 0.20 & setosa \\ 
  2 & 4.90 & 3.00 & 1.40 & 0.20 & setosa \\ 
  3 & 4.70 & 3.20 & 1.30 & 0.20 & setosa \\ 
  4 & 4.60 & 3.10 & 1.50 & 0.20 & setosa \\ 
  5 & 5.00 & 3.60 & 1.40 & 0.20 & setosa \\ 
  6 & 5.40 & 3.90 & 1.70 & 0.40 & setosa \\ 
   \hline
\end{tabular}
\end{table}


}
\end{frame}

%---------------------------------------------------------------------------

\begin{frame}[fragile]
\frametitle{Creating LaTeX Tables with xtable}
\begin{knitrout}\footnotesize
\definecolor{shadecolor}{rgb}{1, 1, 1}\color{fgcolor}\begin{kframe}
\begin{alltt}
\hlopt{?}\hlstd{xtable}
\hlopt{?}\hlstd{print.xtable}
\hlkwd{print}\hlstd{(}\hlkwd{xtable}\hlstd{(}\hlkwd{head}\hlstd{(iris),} \hlkwc{caption}\hlstd{=}\hlstr{"Iris dataset included with R"}\hlstd{),}
      \hlkwc{include.rownames}\hlstd{=}\hlnum{FALSE}\hlstd{,} \hlkwc{size}\hlstd{=}\hlstr{"footnotesize"}\hlstd{)}
\end{alltt}
\end{kframe}
\end{knitrout}


% latex table generated in R 3.0.2 by xtable 1.7-1 package
% Wed Mar  5 15:14:09 2014
\begin{table}[ht]
\centering
{\footnotesize
\begin{tabular}{rrrrl}
  \hline
Sepal.Length & Sepal.Width & Petal.Length & Petal.Width & Species \\ 
  \hline
5.10 & 3.50 & 1.40 & 0.20 & setosa \\ 
  4.90 & 3.00 & 1.40 & 0.20 & setosa \\ 
  4.70 & 3.20 & 1.30 & 0.20 & setosa \\ 
  4.60 & 3.10 & 1.50 & 0.20 & setosa \\ 
  5.00 & 3.60 & 1.40 & 0.20 & setosa \\ 
  5.40 & 3.90 & 1.70 & 0.40 & setosa \\ 
   \hline
\end{tabular}
}
\caption{Iris dataset included with R} 
\end{table}


\end{frame}

%---------------------------------------------------------------------------

\begin{frame}[fragile]
\frametitle{Your Turn}
Modify \href{../code/02-blank.Rnw}{blank.Rnw} to do the following:
\begin{itemize}
\item Use the \texttt{digits} option in \texttt{xtable} to display only one decimal for the Sepal measurements
\item Change the table caption
\item Set the table reference label to "irisdata" and reference it in a sentence outside the code chunk
\item Remove the row number from the table display
\end{itemize}

\begin{knitrout}\footnotesize
\definecolor{shadecolor}{rgb}{1, 1, 1}\color{fgcolor}\begin{kframe}
\begin{alltt}
\hlkwd{library}\hlstd{(xtable)}
\hlkwd{data}\hlstd{(iris)}
\hlkwd{xtable}\hlstd{(}\hlkwd{head}\hlstd{(iris),} \hlkwc{digits}\hlstd{=}\hlkwd{c}\hlstd{(}\hlnum{0}\hlstd{,} \hlnum{1}\hlstd{,} \hlnum{1}\hlstd{,} \hlnum{2}\hlstd{,} \hlnum{2}\hlstd{,} \hlnum{0}\hlstd{))}
\end{alltt}
\end{kframe}
\end{knitrout}

\end{frame}

%---------------------------------------------------------------------------

\begin{frame}[fragile]
\frametitle{stargazer: Latex and R Models}
\begin{knitrout}\footnotesize
\definecolor{shadecolor}{rgb}{1, 1, 1}\color{fgcolor}\begin{kframe}
\begin{alltt}
\hlstd{x} \hlkwb{<-} \hlkwd{seq}\hlstd{(}\hlnum{0}\hlstd{,} \hlnum{10}\hlstd{,} \hlnum{.5}\hlstd{)}
\hlstd{y} \hlkwb{<-} \hlstd{x}\hlopt{+}\hlkwd{rnorm}\hlstd{(}\hlkwd{length}\hlstd{(x))}
\hlstd{data} \hlkwb{<-} \hlkwd{data.frame}\hlstd{(}\hlkwc{x}\hlstd{=x,} \hlkwc{y}\hlstd{=y)}

\hlstd{model} \hlkwb{<-} \hlkwd{lm}\hlstd{(y}\hlopt{~}\hlstd{x,} \hlkwc{data}\hlstd{=data)}
\hlkwd{library}\hlstd{(stargazer)}
\hlkwd{stargazer}\hlstd{(model)}
\end{alltt}
\end{kframe}
\end{knitrout}

\end{frame}

%---------------------------------------------------------------------------

\begin{frame}[fragile]
\frametitle{stargazer: Latex and R Models}
{\footnotesize

% Table created by stargazer v.5.0 by Marek Hlavac, Harvard University. E-mail: hlavac at fas.harvard.edu
% Date and time: Fri, May 23, 2014 - 4:34:17 PM
\begin{tabular}{@{\extracolsep{5pt}}lc} 
\\[-1.8ex]\hline 
\hline \\[-1.8ex] 
 & \multicolumn{1}{c}{\textit{Dependent variable:}} \\ 
\cline{2-2} 
\\[-1.8ex] & y \\ 
\hline \\[-1.8ex] 
 x & 0.939$^{***}$ \\ 
  & (0.058) \\ 
  & \\ 
 Constant & 0.446 \\ 
  & (0.342) \\ 
  & \\ 
\hline \\[-1.8ex] 
Observations & 21 \\ 
R$^{2}$ & 0.931 \\ 
Adjusted R$^{2}$ & 0.928 \\ 
Residual Std. Error & 0.811 (df = 19) \\ 
F Statistic & 258.100$^{***}$ (df = 1; 19) \\ 
\hline 
\hline \\[-1.8ex] 
\textit{Note:}  & \multicolumn{1}{r}{$^{*}$p$<$0.1; $^{**}$p$<$0.05; $^{***}$p$<$0.01} \\ 
\end{tabular} 


}
\end{frame}

%---------------------------------------------------------------------------

\begin{frame}[fragile]
\frametitle{stargazer: Latex and R Models}
{\tiny
\begin{kframe}
\begin{alltt}
\hlstd{intercept.model} \hlkwb{<-} \hlkwd{lm}\hlstd{(y}\hlopt{~}\hlstd{x,} \hlkwc{data}\hlstd{=data)}
\hlstd{nointercept.model} \hlkwb{<-} \hlkwd{lm}\hlstd{(y}\hlopt{~}\hlnum{0}\hlopt{+}\hlstd{x,} \hlkwc{data}\hlstd{=data)}

\hlkwd{stargazer}\hlstd{(intercept.model, nointercept.model,} \hlkwc{float}\hlstd{=F,} \hlkwc{single.row}\hlstd{=}\hlnum{TRUE}\hlstd{,}
          \hlkwc{font.size}\hlstd{=}\hlstr{"scriptsize"}\hlstd{,} \hlkwc{object.names}\hlstd{=T,} \hlkwc{model.numbers}\hlstd{=F)}
\end{alltt}
\end{kframe}
% Table created by stargazer v.5.0 by Marek Hlavac, Harvard University. E-mail: hlavac at fas.harvard.edu
% Date and time: Fri, May 23, 2014 - 4:34:17 PM
\begingroup 
\scriptsize 
\begin{tabular}{@{\extracolsep{5pt}}lcc} 
\\[-1.8ex]\hline 
\hline \\[-1.8ex] 
 & \multicolumn{2}{c}{\textit{Dependent variable:}} \\ 
\cline{2-3} 
\\[-1.8ex] & \multicolumn{2}{c}{y} \\ 
\\[-1.8ex] & intercept.model & nointercept.model\\ 
\hline \\[-1.8ex] 
 x & 0.939$^{***}$ (0.058) & 1.004$^{***}$ (0.031) \\ 
  Constant & 0.446 (0.342) &  \\ 
 \hline \\[-1.8ex] 
Observations & 21 & 21 \\ 
R$^{2}$ & 0.931 & 0.982 \\ 
Adjusted R$^{2}$ & 0.928 & 0.981 \\ 
Residual Std. Error & 0.811 (df = 19) & 0.825 (df = 20) \\ 
F Statistic & 258.100$^{***}$ (df = 1; 19) & 1,063.000$^{***}$ (df = 1; 20) \\ 
\hline 
\hline \\[-1.8ex] 
\textit{Note:}  & \multicolumn{2}{r}{$^{*}$p$<$0.1; $^{**}$p$<$0.05; $^{***}$p$<$0.01} \\ 
\end{tabular} 
\endgroup 


}
\end{frame}

%---------------------------------------------------------------------------

\begin{frame}[fragile]
\frametitle{Your Turn}
Using the iris data:
\begin{itemize}
\item Model sepal width by petal width; report the results in a LaTeX table using stargazer
\item Add a caption to your table using the \texttt{title=} option
\item Change the independent variable label to read ``Petal Width"\\ Hint: \texttt{dep.var.labels=""}
\item Change the dependent variable label to ``Sepal Width"\\ Hint: \texttt{covariate.labels=""}
\item In the text below the table, report and discuss the correlation between Sepal and Petal Width using \texttt{\textbackslash Sexpr\{\}}.
\end{itemize}

\begin{knitrout}\footnotesize
\definecolor{shadecolor}{rgb}{1, 1, 1}\color{fgcolor}\begin{kframe}
\begin{alltt}
\hlkwd{data}\hlstd{(iris)}
\hlstd{model} \hlkwb{<-} \hlkwd{lm}\hlstd{(iris}\hlopt{$}\hlstd{Sepal.Width}\hlopt{~}\hlstd{iris}\hlopt{$}\hlstd{Petal.Width)}
\hlopt{?}\hlstd{stargazer}
\end{alltt}
\end{kframe}
\end{knitrout}


{\footnotesize Remember to turn messages off!}
\end{frame}

%---------------------------------------------------------------------------

\begin{frame}[fragile]
\frametitle{Hmisc: Latex and R Objects}

Workhorse function: \texttt{latex()}\bigskip

\begin{itemize}
\item Very flexible - handles S3 and S4 classes, lists, data frames, matrices...
\item Lots of options for formatting
\item Has a preview function
\item Other packages are easier to use
\end{itemize}
\end{frame}

%---------------------------------------------------------------------------

\begin{frame}[fragile]
\frametitle{Hmisc: Latex and R Objects}
\begin{kframe}
\begin{alltt}
\hlkwd{library}\hlstd{(Hmisc)}
\hlstd{my.mean} \hlkwb{<-} \hlkwa{function}\hlstd{(}\hlkwc{x}\hlstd{)\{}
  \hlkwd{sum}\hlstd{(x,} \hlkwc{na.rm}\hlstd{=}\hlnum{TRUE}\hlstd{)}\hlopt{/}\hlkwd{length}\hlstd{(x)}
\hlstd{\}}
\hlkwd{latex}\hlstd{(my.mean,} \hlkwc{file}\hlstd{=}\hlstr{""}\hlstd{)}
\end{alltt}
\end{kframe}\begin{alltt}
my.mean \(\leftarrow\) function (x) 
\{
    sum(x, na.rm = TRUE)/length(x)
\}
\end{alltt}



\end{frame}

%---------------------------------------------------------------------------

\begin{frame}[fragile]
\frametitle{Hmisc: Latex and R Objects}
Functions can also be done with knitr directly: 
\begin{kframe}
\begin{alltt}
\hlstd{my.mean}
\end{alltt}
\end{kframe}function(x){
  sum(x, na.rm=TRUE)/length(x)
}



\end{frame}

%---------------------------------------------------------------------------

\begin{frame}[fragile]
\frametitle{Hmisc: Latex and R Objects}
Output specific model information with \texttt{latex()}
\begin{kframe}
\begin{alltt}
\hlstd{model} \hlkwb{<-} \hlkwd{lm}\hlstd{(y}\hlopt{~}\hlstd{x,} \hlkwc{data}\hlstd{=data)}
\hlstd{model.tab} \hlkwb{<-} \hlkwd{cbind}\hlstd{(}\hlkwc{terms}\hlstd{=}\hlkwd{c}\hlstd{(}\hlstr{"Intercept"}\hlstd{,} \hlstr{"x"}\hlstd{),}
                   \hlkwc{Estimates}\hlstd{=}\hlkwd{round}\hlstd{(model}\hlopt{$}\hlstd{coefficients,} \hlnum{3}\hlstd{),}
                   \hlkwc{SE}\hlstd{=}\hlkwd{round}\hlstd{(}\hlkwd{sqrt}\hlstd{(}\hlkwd{diag}\hlstd{(}\hlkwd{vcov}\hlstd{(model))),} \hlnum{3}\hlstd{))}
\hlkwd{rownames}\hlstd{(model.tab)} \hlkwb{<-} \hlkwa{NULL}
\hlkwd{latex}\hlstd{(model.tab,} \hlkwc{file}\hlstd{=}\hlstr{""}\hlstd{)}
\end{alltt}
\end{kframe}%latex.default(model.tab, file = "")%
\begin{table}[!tbp]
\begin{center}
\begin{tabular}{lll}
\hline\hline
\multicolumn{1}{c}{terms}&\multicolumn{1}{c}{Estimates}&\multicolumn{1}{c}{SE}\tabularnewline
\hline
Intercept&0.446&0.342\tabularnewline
x&0.939&0.058\tabularnewline
\hline
\end{tabular}\end{center}\end{table}



\end{frame}

%---------------------------------------------------------------------------


\end{document}
